% =========================================================
% Slides Beamer Metropolis 16:9 — Chapter 2 (Decomposition)
% Time Series Analysis — Aymen Ben Brik (aymen.benbrik@esprit.tn)
%
% Merged from _archives/source-slides/Time series Decomposition 1.tex
% (Trend, 414 lines, Madrid) and Time series decomposition 2.tex
% (Seasonality, 293 lines, Madrid).
% Compile : pdflatex slides.tex (twice).
% =========================================================
\documentclass[aspectratio=169]{beamer}

% =========================================================
% Beamer preamble — Time Series Analysis module
% Theme : Metropolis, aspect ratio 16:9
% Author : Aymen Ben Brik (aymen.benbrik@esprit.tn)
% Esprit School of Business
%
% Include from each chapterX/slides.tex via :
%   % =========================================================
% Beamer preamble — Time Series Analysis module
% Theme : Metropolis, aspect ratio 16:9
% Author : Aymen Ben Brik (aymen.benbrik@esprit.tn)
% Esprit School of Business
%
% Include from each chapterX/slides.tex via :
%   % =========================================================
% Beamer preamble — Time Series Analysis module
% Theme : Metropolis, aspect ratio 16:9
% Author : Aymen Ben Brik (aymen.benbrik@esprit.tn)
% Esprit School of Business
%
% Include from each chapterX/slides.tex via :
%   \input{../preamble/slides-preamble.tex}
% AFTER the line \documentclass[aspectratio=169]{beamer}.
% =========================================================

\usepackage[utf8]{inputenc}
\usepackage[T1]{fontenc}
\usepackage[english]{babel}
\usepackage{lmodern}
\usepackage{amsmath, amssymb, mathtools, bm}
\usepackage{graphicx}
\usepackage{booktabs}
\usepackage{array}
\usepackage{multirow}
\usepackage{colortbl}
\usepackage{xcolor}
\usepackage{tikz}
\usetikzlibrary{positioning, arrows.meta, calc, shapes, fit, backgrounds, matrix, decorations.pathreplacing}
\usepackage{listings}
\usepackage[most]{tcolorbox}

% --- Theme ---
\usetheme{metropolis}
\setbeamertemplate{navigation symbols}{}
\metroset{block=fill}

% --- Esprit colour palette ---
\definecolor{espritBlue}{RGB}{30, 60, 110}
\definecolor{espritAccent}{RGB}{200, 80, 30}
\setbeamercolor{progress bar}{fg=espritAccent}
\setbeamercolor{title separator}{fg=espritAccent}
\setbeamercolor{frametitle}{bg=espritBlue}

% --- Code listings (Python + R, with UTF-8 literate) ---
\definecolor{codebg}{RGB}{248,248,248}
\definecolor{codeframe}{RGB}{220,220,220}
\definecolor{keyword}{RGB}{0,82,155}
\definecolor{comment}{RGB}{88,110,117}
\definecolor{string}{RGB}{133,153,0}

\lstdefinestyle{pythonstyle}{
  language=Python,
  basicstyle=\ttfamily\scriptsize,
  keywordstyle=\color{keyword}\bfseries,
  commentstyle=\color{comment}\itshape,
  stringstyle=\color{string},
  backgroundcolor=\color{codebg},
  frame=single,
  rulecolor=\color{codeframe},
  frameround=tttt,
  breaklines=true,
  showstringspaces=false,
  tabsize=2,
  literate=
    {é}{{\'e}}1 {è}{{\`e}}1 {ê}{{\^e}}1 {à}{{\`a}}1 {â}{{\^a}}1
    {ç}{{\c{c}}}1 {²}{{$^2$}}1 {°}{{${}^\circ$}}1
    {α}{{$\alpha$}}1 {β}{{$\beta$}}1 {γ}{{$\gamma$}}1
    {δ}{{$\delta$}}1 {θ}{{$\theta$}}1 {λ}{{$\lambda$}}1
    {μ}{{$\mu$}}1 {σ}{{$\sigma$}}1 {ε}{{$\varepsilon$}}1
    {ρ}{{$\rho$}}1 {τ}{{$\tau$}}1 {φ}{{$\phi$}}1
    {≤}{{$\leq$}}1 {≥}{{$\geq$}}1 {≈}{{$\approx$}}1
}
\lstdefinestyle{Rstyle}{
  language=R,
  basicstyle=\ttfamily\scriptsize,
  keywordstyle=\color{keyword}\bfseries,
  commentstyle=\color{comment}\itshape,
  stringstyle=\color{string},
  backgroundcolor=\color{codebg},
  frame=single,
  rulecolor=\color{codeframe},
  frameround=tttt,
  breaklines=true,
  showstringspaces=false,
  tabsize=2
}
\lstset{style=pythonstyle}

% --- Common macros ---
\input{../preamble/macros.tex}

% --- tcolorbox boxes for slides ---
\definecolor{boxTh}{RGB}{120, 30, 30}
\definecolor{boxProp}{RGB}{30, 100, 60}
\definecolor{boxDef}{RGB}{30, 60, 110}
\definecolor{boxEx}{RGB}{180, 120, 0}

\newtcolorbox{infobox}[1][]{enhanced, colback=espritBlue!5, colframe=espritBlue,
  fonttitle=\bfseries\small,
  title={\ifstrempty{#1}{}{#1}},
  boxrule=0.4pt, arc=2pt, left=4pt, right=4pt, top=2pt, bottom=2pt}
\newtcolorbox{warnbox}[1][]{enhanced, colback=espritAccent!5, colframe=espritAccent,
  fonttitle=\bfseries\small,
  title={\ifstrempty{#1}{}{#1}},
  boxrule=0.4pt, arc=2pt, left=4pt, right=4pt, top=2pt, bottom=2pt}
\newtcolorbox{thbox}[1][]{enhanced, colback=boxTh!5, colframe=boxTh,
  fonttitle=\bfseries\small,
  title={Theorem\ifstrempty{#1}{}{ -- #1}},
  boxrule=0.5pt, arc=2pt, left=4pt, right=4pt, top=2pt, bottom=2pt}
\newtcolorbox{propbox}[1][]{enhanced, colback=boxProp!5, colframe=boxProp,
  fonttitle=\bfseries\small,
  title={Proposition\ifstrempty{#1}{}{ -- #1}},
  boxrule=0.5pt, arc=2pt, left=4pt, right=4pt, top=2pt, bottom=2pt}
\newtcolorbox{defbox}[1][]{enhanced, colback=boxDef!5, colframe=boxDef,
  fonttitle=\bfseries\small,
  title={Definition\ifstrempty{#1}{}{ -- #1}},
  boxrule=0.5pt, arc=2pt, left=4pt, right=4pt, top=2pt, bottom=2pt}
\newtcolorbox{exbox}[1][]{enhanced, colback=boxEx!5, colframe=boxEx,
  fonttitle=\bfseries\small,
  title={Example\ifstrempty{#1}{}{ -- #1}},
  boxrule=0.5pt, arc=2pt, left=4pt, right=4pt, top=2pt, bottom=2pt}

% AFTER the line \documentclass[aspectratio=169]{beamer}.
% =========================================================

\usepackage[utf8]{inputenc}
\usepackage[T1]{fontenc}
\usepackage[english]{babel}
\usepackage{lmodern}
\usepackage{amsmath, amssymb, mathtools, bm}
\usepackage{graphicx}
\usepackage{booktabs}
\usepackage{array}
\usepackage{multirow}
\usepackage{colortbl}
\usepackage{xcolor}
\usepackage{tikz}
\usetikzlibrary{positioning, arrows.meta, calc, shapes, fit, backgrounds, matrix, decorations.pathreplacing}
\usepackage{listings}
\usepackage[most]{tcolorbox}

% --- Theme ---
\usetheme{metropolis}
\setbeamertemplate{navigation symbols}{}
\metroset{block=fill}

% --- Esprit colour palette ---
\definecolor{espritBlue}{RGB}{30, 60, 110}
\definecolor{espritAccent}{RGB}{200, 80, 30}
\setbeamercolor{progress bar}{fg=espritAccent}
\setbeamercolor{title separator}{fg=espritAccent}
\setbeamercolor{frametitle}{bg=espritBlue}

% --- Code listings (Python + R, with UTF-8 literate) ---
\definecolor{codebg}{RGB}{248,248,248}
\definecolor{codeframe}{RGB}{220,220,220}
\definecolor{keyword}{RGB}{0,82,155}
\definecolor{comment}{RGB}{88,110,117}
\definecolor{string}{RGB}{133,153,0}

\lstdefinestyle{pythonstyle}{
  language=Python,
  basicstyle=\ttfamily\scriptsize,
  keywordstyle=\color{keyword}\bfseries,
  commentstyle=\color{comment}\itshape,
  stringstyle=\color{string},
  backgroundcolor=\color{codebg},
  frame=single,
  rulecolor=\color{codeframe},
  frameround=tttt,
  breaklines=true,
  showstringspaces=false,
  tabsize=2,
  literate=
    {é}{{\'e}}1 {è}{{\`e}}1 {ê}{{\^e}}1 {à}{{\`a}}1 {â}{{\^a}}1
    {ç}{{\c{c}}}1 {²}{{$^2$}}1 {°}{{${}^\circ$}}1
    {α}{{$\alpha$}}1 {β}{{$\beta$}}1 {γ}{{$\gamma$}}1
    {δ}{{$\delta$}}1 {θ}{{$\theta$}}1 {λ}{{$\lambda$}}1
    {μ}{{$\mu$}}1 {σ}{{$\sigma$}}1 {ε}{{$\varepsilon$}}1
    {ρ}{{$\rho$}}1 {τ}{{$\tau$}}1 {φ}{{$\phi$}}1
    {≤}{{$\leq$}}1 {≥}{{$\geq$}}1 {≈}{{$\approx$}}1
}
\lstdefinestyle{Rstyle}{
  language=R,
  basicstyle=\ttfamily\scriptsize,
  keywordstyle=\color{keyword}\bfseries,
  commentstyle=\color{comment}\itshape,
  stringstyle=\color{string},
  backgroundcolor=\color{codebg},
  frame=single,
  rulecolor=\color{codeframe},
  frameround=tttt,
  breaklines=true,
  showstringspaces=false,
  tabsize=2
}
\lstset{style=pythonstyle}

% --- Common macros ---
% =========================================================
% Common math macros — Time Series Analysis module
% Author : Aymen Ben Brik (aymen.benbrik@esprit.tn)
% Esprit School of Business
% =========================================================

% --- Standard sets ---
\providecommand{\R}{\mathbb{R}}
\providecommand{\N}{\mathbb{N}}
\providecommand{\Z}{\mathbb{Z}}
\providecommand{\Q}{\mathbb{Q}}
\providecommand{\C}{\mathbb{C}}

% --- Probability / statistics ---
\providecommand{\E}{\mathbb{E}}
\providecommand{\Prob}{\mathbb{P}}
\providecommand{\Var}{\mathrm{Var}}
\providecommand{\Cov}{\mathrm{Cov}}
\providecommand{\Corr}{\mathrm{Corr}}
\providecommand{\indep}{\perp\!\!\!\perp}
\providecommand{\iid}{\overset{\text{i.i.d.}}{\sim}}
\providecommand{\given}{\,\big\vert\,}

% --- Estimators / hat ---
\providecommand{\hatm}{\widehat{m}}
\providecommand{\hatsig}{\widehat{\sigma}}
\providecommand{\hatrho}{\widehat{\rho}}
\providecommand{\hatphi}{\widehat{\phi}}
\providecommand{\hattheta}{\widehat{\theta}}

% --- Time series notation ---
\providecommand{\Lop}{L}                      % lag operator
\providecommand{\Diff}{\Delta}                % difference operator
\providecommand{\AR}[1]{\mathrm{AR}(#1)}
\providecommand{\MAm}[1]{\mathrm{MA}(#1)}
\providecommand{\ARMA}[2]{\mathrm{ARMA}(#1,#2)}
\providecommand{\ARIMA}[3]{\mathrm{ARIMA}(#1,#2,#3)}
\providecommand{\SARIMA}[6]{\mathrm{SARIMA}(#1,#2,#3)(#4,#5,#6)}
\providecommand{\ARCH}[1]{\mathrm{ARCH}(#1)}
\providecommand{\GARCH}[2]{\mathrm{GARCH}(#1,#2)}
\providecommand{\WN}{\mathrm{WN}}              % white noise
\providecommand{\MSE}{\mathrm{MSE}}
\providecommand{\MAE}{\mathrm{MAE}}
\providecommand{\RMSE}{\mathrm{RMSE}}
\providecommand{\AIC}{\mathrm{AIC}}
\providecommand{\BIC}{\mathrm{BIC}}
\providecommand{\ACF}{\mathrm{ACF}}
\providecommand{\PACF}{\mathrm{PACF}}

% --- Linear algebra ---
\providecommand{\transp}[1]{#1^{\!\top}}
\providecommand{\norm}[1]{\left\lVert #1 \right\rVert}
\providecommand{\abs}[1]{\left\lvert #1 \right\rvert}
\providecommand{\inner}[2]{\langle #1, #2 \rangle}
\providecommand{\diag}{\mathrm{diag}}
\providecommand{\tr}{\mathrm{tr}}

% --- Bold vectors / matrices ---
\providecommand{\vx}{\bm{x}}
\providecommand{\vy}{\bm{y}}
\providecommand{\vz}{\bm{z}}
\providecommand{\vh}{\bm{h}}
\providecommand{\vu}{\bm{u}}
\providecommand{\vv}{\bm{v}}
\providecommand{\vw}{\bm{w}}
\providecommand{\vb}{\bm{b}}
\providecommand{\vW}{\bm{W}}
\providecommand{\vSigma}{\bm{\Sigma}}
\providecommand{\veps}{\bm{\varepsilon}}

% --- Author identity (reused on titlepages) ---
\providecommand{\auteurNom}{Aymen Ben Brik}
\providecommand{\auteurEmail}{aymen.benbrik@esprit.tn}
\providecommand{\auteurInstitution}{Esprit School of Business}
\providecommand{\auteurRole}{Module head}
\providecommand{\moduleNom}{Time Series Analysis}
\providecommand{\anneeUniversitaire}{2025--2026}


% --- tcolorbox boxes for slides ---
\definecolor{boxTh}{RGB}{120, 30, 30}
\definecolor{boxProp}{RGB}{30, 100, 60}
\definecolor{boxDef}{RGB}{30, 60, 110}
\definecolor{boxEx}{RGB}{180, 120, 0}

\newtcolorbox{infobox}[1][]{enhanced, colback=espritBlue!5, colframe=espritBlue,
  fonttitle=\bfseries\small,
  title={\ifstrempty{#1}{}{#1}},
  boxrule=0.4pt, arc=2pt, left=4pt, right=4pt, top=2pt, bottom=2pt}
\newtcolorbox{warnbox}[1][]{enhanced, colback=espritAccent!5, colframe=espritAccent,
  fonttitle=\bfseries\small,
  title={\ifstrempty{#1}{}{#1}},
  boxrule=0.4pt, arc=2pt, left=4pt, right=4pt, top=2pt, bottom=2pt}
\newtcolorbox{thbox}[1][]{enhanced, colback=boxTh!5, colframe=boxTh,
  fonttitle=\bfseries\small,
  title={Theorem\ifstrempty{#1}{}{ -- #1}},
  boxrule=0.5pt, arc=2pt, left=4pt, right=4pt, top=2pt, bottom=2pt}
\newtcolorbox{propbox}[1][]{enhanced, colback=boxProp!5, colframe=boxProp,
  fonttitle=\bfseries\small,
  title={Proposition\ifstrempty{#1}{}{ -- #1}},
  boxrule=0.5pt, arc=2pt, left=4pt, right=4pt, top=2pt, bottom=2pt}
\newtcolorbox{defbox}[1][]{enhanced, colback=boxDef!5, colframe=boxDef,
  fonttitle=\bfseries\small,
  title={Definition\ifstrempty{#1}{}{ -- #1}},
  boxrule=0.5pt, arc=2pt, left=4pt, right=4pt, top=2pt, bottom=2pt}
\newtcolorbox{exbox}[1][]{enhanced, colback=boxEx!5, colframe=boxEx,
  fonttitle=\bfseries\small,
  title={Example\ifstrempty{#1}{}{ -- #1}},
  boxrule=0.5pt, arc=2pt, left=4pt, right=4pt, top=2pt, bottom=2pt}

% AFTER the line \documentclass[aspectratio=169]{beamer}.
% =========================================================

\usepackage[utf8]{inputenc}
\usepackage[T1]{fontenc}
\usepackage[english]{babel}
\usepackage{lmodern}
\usepackage{amsmath, amssymb, mathtools, bm}
\usepackage{graphicx}
\usepackage{booktabs}
\usepackage{array}
\usepackage{multirow}
\usepackage{colortbl}
\usepackage{xcolor}
\usepackage{tikz}
\usetikzlibrary{positioning, arrows.meta, calc, shapes, fit, backgrounds, matrix, decorations.pathreplacing}
\usepackage{listings}
\usepackage[most]{tcolorbox}

% --- Theme ---
\usetheme{metropolis}
\setbeamertemplate{navigation symbols}{}
\metroset{block=fill}

% --- Esprit colour palette ---
\definecolor{espritBlue}{RGB}{30, 60, 110}
\definecolor{espritAccent}{RGB}{200, 80, 30}
\setbeamercolor{progress bar}{fg=espritAccent}
\setbeamercolor{title separator}{fg=espritAccent}
\setbeamercolor{frametitle}{bg=espritBlue}

% --- Code listings (Python + R, with UTF-8 literate) ---
\definecolor{codebg}{RGB}{248,248,248}
\definecolor{codeframe}{RGB}{220,220,220}
\definecolor{keyword}{RGB}{0,82,155}
\definecolor{comment}{RGB}{88,110,117}
\definecolor{string}{RGB}{133,153,0}

\lstdefinestyle{pythonstyle}{
  language=Python,
  basicstyle=\ttfamily\scriptsize,
  keywordstyle=\color{keyword}\bfseries,
  commentstyle=\color{comment}\itshape,
  stringstyle=\color{string},
  backgroundcolor=\color{codebg},
  frame=single,
  rulecolor=\color{codeframe},
  frameround=tttt,
  breaklines=true,
  showstringspaces=false,
  tabsize=2,
  literate=
    {é}{{\'e}}1 {è}{{\`e}}1 {ê}{{\^e}}1 {à}{{\`a}}1 {â}{{\^a}}1
    {ç}{{\c{c}}}1 {²}{{$^2$}}1 {°}{{${}^\circ$}}1
    {α}{{$\alpha$}}1 {β}{{$\beta$}}1 {γ}{{$\gamma$}}1
    {δ}{{$\delta$}}1 {θ}{{$\theta$}}1 {λ}{{$\lambda$}}1
    {μ}{{$\mu$}}1 {σ}{{$\sigma$}}1 {ε}{{$\varepsilon$}}1
    {ρ}{{$\rho$}}1 {τ}{{$\tau$}}1 {φ}{{$\phi$}}1
    {≤}{{$\leq$}}1 {≥}{{$\geq$}}1 {≈}{{$\approx$}}1
}
\lstdefinestyle{Rstyle}{
  language=R,
  basicstyle=\ttfamily\scriptsize,
  keywordstyle=\color{keyword}\bfseries,
  commentstyle=\color{comment}\itshape,
  stringstyle=\color{string},
  backgroundcolor=\color{codebg},
  frame=single,
  rulecolor=\color{codeframe},
  frameround=tttt,
  breaklines=true,
  showstringspaces=false,
  tabsize=2
}
\lstset{style=pythonstyle}

% --- Common macros ---
% =========================================================
% Common math macros — Time Series Analysis module
% Author : Aymen Ben Brik (aymen.benbrik@esprit.tn)
% Esprit School of Business
% =========================================================

% --- Standard sets ---
\providecommand{\R}{\mathbb{R}}
\providecommand{\N}{\mathbb{N}}
\providecommand{\Z}{\mathbb{Z}}
\providecommand{\Q}{\mathbb{Q}}
\providecommand{\C}{\mathbb{C}}

% --- Probability / statistics ---
\providecommand{\E}{\mathbb{E}}
\providecommand{\Prob}{\mathbb{P}}
\providecommand{\Var}{\mathrm{Var}}
\providecommand{\Cov}{\mathrm{Cov}}
\providecommand{\Corr}{\mathrm{Corr}}
\providecommand{\indep}{\perp\!\!\!\perp}
\providecommand{\iid}{\overset{\text{i.i.d.}}{\sim}}
\providecommand{\given}{\,\big\vert\,}

% --- Estimators / hat ---
\providecommand{\hatm}{\widehat{m}}
\providecommand{\hatsig}{\widehat{\sigma}}
\providecommand{\hatrho}{\widehat{\rho}}
\providecommand{\hatphi}{\widehat{\phi}}
\providecommand{\hattheta}{\widehat{\theta}}

% --- Time series notation ---
\providecommand{\Lop}{L}                      % lag operator
\providecommand{\Diff}{\Delta}                % difference operator
\providecommand{\AR}[1]{\mathrm{AR}(#1)}
\providecommand{\MAm}[1]{\mathrm{MA}(#1)}
\providecommand{\ARMA}[2]{\mathrm{ARMA}(#1,#2)}
\providecommand{\ARIMA}[3]{\mathrm{ARIMA}(#1,#2,#3)}
\providecommand{\SARIMA}[6]{\mathrm{SARIMA}(#1,#2,#3)(#4,#5,#6)}
\providecommand{\ARCH}[1]{\mathrm{ARCH}(#1)}
\providecommand{\GARCH}[2]{\mathrm{GARCH}(#1,#2)}
\providecommand{\WN}{\mathrm{WN}}              % white noise
\providecommand{\MSE}{\mathrm{MSE}}
\providecommand{\MAE}{\mathrm{MAE}}
\providecommand{\RMSE}{\mathrm{RMSE}}
\providecommand{\AIC}{\mathrm{AIC}}
\providecommand{\BIC}{\mathrm{BIC}}
\providecommand{\ACF}{\mathrm{ACF}}
\providecommand{\PACF}{\mathrm{PACF}}

% --- Linear algebra ---
\providecommand{\transp}[1]{#1^{\!\top}}
\providecommand{\norm}[1]{\left\lVert #1 \right\rVert}
\providecommand{\abs}[1]{\left\lvert #1 \right\rvert}
\providecommand{\inner}[2]{\langle #1, #2 \rangle}
\providecommand{\diag}{\mathrm{diag}}
\providecommand{\tr}{\mathrm{tr}}

% --- Bold vectors / matrices ---
\providecommand{\vx}{\bm{x}}
\providecommand{\vy}{\bm{y}}
\providecommand{\vz}{\bm{z}}
\providecommand{\vh}{\bm{h}}
\providecommand{\vu}{\bm{u}}
\providecommand{\vv}{\bm{v}}
\providecommand{\vw}{\bm{w}}
\providecommand{\vb}{\bm{b}}
\providecommand{\vW}{\bm{W}}
\providecommand{\vSigma}{\bm{\Sigma}}
\providecommand{\veps}{\bm{\varepsilon}}

% --- Author identity (reused on titlepages) ---
\providecommand{\auteurNom}{Aymen Ben Brik}
\providecommand{\auteurEmail}{aymen.benbrik@esprit.tn}
\providecommand{\auteurInstitution}{Esprit School of Business}
\providecommand{\auteurRole}{Module head}
\providecommand{\moduleNom}{Time Series Analysis}
\providecommand{\anneeUniversitaire}{2025--2026}


% --- tcolorbox boxes for slides ---
\definecolor{boxTh}{RGB}{120, 30, 30}
\definecolor{boxProp}{RGB}{30, 100, 60}
\definecolor{boxDef}{RGB}{30, 60, 110}
\definecolor{boxEx}{RGB}{180, 120, 0}

\newtcolorbox{infobox}[1][]{enhanced, colback=espritBlue!5, colframe=espritBlue,
  fonttitle=\bfseries\small,
  title={\ifstrempty{#1}{}{#1}},
  boxrule=0.4pt, arc=2pt, left=4pt, right=4pt, top=2pt, bottom=2pt}
\newtcolorbox{warnbox}[1][]{enhanced, colback=espritAccent!5, colframe=espritAccent,
  fonttitle=\bfseries\small,
  title={\ifstrempty{#1}{}{#1}},
  boxrule=0.4pt, arc=2pt, left=4pt, right=4pt, top=2pt, bottom=2pt}
\newtcolorbox{thbox}[1][]{enhanced, colback=boxTh!5, colframe=boxTh,
  fonttitle=\bfseries\small,
  title={Theorem\ifstrempty{#1}{}{ -- #1}},
  boxrule=0.5pt, arc=2pt, left=4pt, right=4pt, top=2pt, bottom=2pt}
\newtcolorbox{propbox}[1][]{enhanced, colback=boxProp!5, colframe=boxProp,
  fonttitle=\bfseries\small,
  title={Proposition\ifstrempty{#1}{}{ -- #1}},
  boxrule=0.5pt, arc=2pt, left=4pt, right=4pt, top=2pt, bottom=2pt}
\newtcolorbox{defbox}[1][]{enhanced, colback=boxDef!5, colframe=boxDef,
  fonttitle=\bfseries\small,
  title={Definition\ifstrempty{#1}{}{ -- #1}},
  boxrule=0.5pt, arc=2pt, left=4pt, right=4pt, top=2pt, bottom=2pt}
\newtcolorbox{exbox}[1][]{enhanced, colback=boxEx!5, colframe=boxEx,
  fonttitle=\bfseries\small,
  title={Example\ifstrempty{#1}{}{ -- #1}},
  boxrule=0.5pt, arc=2pt, left=4pt, right=4pt, top=2pt, bottom=2pt}


\title[\moduleNom\ -- Ch.2]{Time Series Decomposition}
\subtitle{Trend $\cdot$ Seasonality $\cdot$ Residuals}
\author{\auteurNom \\ \footnotesize\auteurEmail}
\institute{\auteurInstitution}
\date{\anneeUniversitaire}

\begin{document}

\begin{frame}\titlepage\end{frame}

\begin{frame}{Learning objectives}
  By the end of this chapter, you will be able to:
  \begin{itemize}
    \item write down the additive and multiplicative decomposition
          models for a time series;
    \item estimate the \textbf{trend} by moving average, parametric
          (polynomial) and non-parametric (LOESS) regression;
    \item estimate the \textbf{seasonality} by seasonal averaging,
          dummy regression, and cosine--sine (Fourier) models;
    \item produce a deseasonalised, detrended residual and check it
          for stationarity (preview of Chapter~3).
  \end{itemize}
\end{frame}

\begin{frame}{Chapter outline}
  \begin{enumerate}
    \item The decomposition model (additive vs multiplicative)
    \item Trend estimation (moving average, parametric, LOESS)
    \item Seasonality estimation (averaging, dummy, Fourier)
    \item Residuals and diagnostics
    \item Full decomposition workflow
  \end{enumerate}
\end{frame}

% =====================================================================
\section{The decomposition model}

\begin{frame}{1.~Additive and multiplicative models}
  \begin{defbox}[Two classical decompositions]
    \[
      \text{Additive:}\qquad Y_t = T_t + S_t + \varepsilon_t,
    \]
    \[
      \text{Multiplicative:}\qquad Y_t = T_t \cdot S_t \cdot \varepsilon_t.
    \]
    Here $T_t$ is the trend, $S_t$ the seasonality (period $d$), and
    $\varepsilon_t$ a stationary residual.
  \end{defbox}
  \medskip
  \begin{itemize}
    \item Additive model: seasonal amplitude is roughly \emph{constant}
          over time.
    \item Multiplicative model: amplitude \emph{grows with the level}
          ($Y_t > 0$); take $\log Y_t$ to recover the additive form.
    \item Goal: estimate $\widehat{T}_t$ and $\widehat{S}_t$, then
          $\widehat{\varepsilon}_t = Y_t - \widehat{T}_t - \widehat{S}_t$
          should look stationary.
  \end{itemize}
\end{frame}

% =====================================================================
\section{Trend estimation}

\begin{frame}{2.~Symmetric (centred) moving average}
  \begin{defbox}[Moving average]
    For odd window $2k+1$ (with $k$ small relative to $T$):
    \[
      \widehat{T}_t \;=\; \frac{1}{2k+1}\sum_{i=-k}^{k} Y_{t+i},
      \quad t = k+1, \dots, T-k.
    \]
  \end{defbox}
  \medskip
  \textbf{Properties.}
  \begin{itemize}
    \item \textbf{Preserves linear trends:} if $Y_t = a + b\,t$, then
          $\widehat{T}_t = a + b\,t$ (exact).
    \item \textbf{Removes additive periodic seasonality with period $d = 2k+1$.}
    \item Loses $k$ observations at each end (boundary effect).
  \end{itemize}
\end{frame}

\begin{frame}{2.~Numerical example: $k = 1$ (3-point MA)}
  Series $Y = (10, 12, 15, 14, 16, 18, 20)$:
  \[
  \begin{array}{rl}
    \widehat{T}_2 &= \tfrac{10 + 12 + 15}{3} \approx 12.33,\\
    \widehat{T}_3 &= \tfrac{12 + 15 + 14}{3} \approx 13.67,\\
    \widehat{T}_4 &= \tfrac{15 + 14 + 16}{3} = 15.00,\\
    \widehat{T}_5 &= \tfrac{14 + 16 + 18}{3} = 16.00,\\
    \widehat{T}_6 &= \tfrac{16 + 18 + 20}{3} = 18.00.
  \end{array}
  \]
  Boundaries $t = 1$ and $t = 7$ are not estimated ($k = 1$ data
  point lost on each side).
\end{frame}

\begin{frame}{2.~Parametric trend (polynomial regression)}
  \begin{defbox}[Polynomial trend]
    \[
      T_t \;=\; \beta_0 + \beta_1\,t + \beta_2\,t^2 + \dots + \beta_p\,t^p.
    \]
  \end{defbox}
  \medskip
  \begin{itemize}
    \item $p = 1$ : straight line. Estimate $(\beta_0, \beta_1)$ by OLS.
    \item $p = 2$ : quadratic. Captures curvature.
    \item $p$ large : risk of overfitting and bad extrapolation.
  \end{itemize}
  \medskip
  \begin{infobox}
    Choose $p$ via AIC / BIC, cross-validation or domain knowledge.
    Refrain from polynomials of order $> 3$ — they are unstable at
    boundaries.
  \end{infobox}
\end{frame}

\begin{frame}[fragile]{2.~Linear trend in Python (OLS)}
\begin{lstlisting}
import numpy as np
from sklearn.linear_model import LinearRegression

t = np.arange(1, len(y) + 1).reshape(-1, 1)
model = LinearRegression().fit(t, y)
trend = model.predict(t)
beta0, beta1 = model.intercept_, model.coef_[0]
\end{lstlisting}
  \medskip
  \textbf{Reading.} $\widehat\beta_1$ is the per-period slope of the trend.
  Subtract \texttt{trend} from \texttt{y} to detrend.
\end{frame}

\begin{frame}{2.~Non-parametric trend (LOESS / splines)}
  \begin{defbox}[Local polynomial / LOESS]
    For each $t$, fit a low-degree polynomial \emph{locally} on a
    window around $t$ with weights decreasing in $|t' - t|$ (typically
    a tricube kernel). The window width is a hyper-parameter (\emph{span}).
  \end{defbox}
  \medskip
  \begin{itemize}
    \item Flexibility: handles non-linear, non-parametric trends.
    \item Sensitive to boundary effects (less than MA but still).
    \item Other choices: \textbf{splines} (penalised cubic), \textbf{wavelets}.
  \end{itemize}
  \medskip
  \begin{warnbox}
    Hyper-parameter trade-off: small span $\Rightarrow$ overfitting,
    large span $\Rightarrow$ over-smoothing.
  \end{warnbox}
\end{frame}

\begin{frame}[fragile]{2.~LOESS in Python}
\begin{lstlisting}
from statsmodels.nonparametric.smoothers_lowess import lowess

trend = lowess(y, t, frac=0.2, return_sorted=False)

# detrended series
detrended = y - trend
\end{lstlisting}
  \medskip
  \texttt{frac} = fraction of data in each local window
  ($\sim$ smoothing parameter).
\end{frame}

% =====================================================================
\section{Seasonality estimation}

\begin{frame}{3.~Seasonal averaging (after detrending)}
  Once the trend is removed, observations at the same season $k$
  ($k = 1, \dots, d$) are averaged:
  \[
    \widehat{S}_k \;=\; \frac{1}{n_k}\sum_{t \in \text{season } k}
                    \bigl(Y_t - \widehat{T}_t\bigr).
  \]
  \medskip
  \begin{itemize}
    \item For monthly data: $d = 12$, one $\widehat{S}_k$ per calendar
          month.
    \item Convention: re-centre so $\sum_k \widehat{S}_k = 0$
          (additive case).
    \item Deseasonalised series:
          $Y_t^{\text{adj}} = Y_t - \widehat{S}_{k(t)}$.
  \end{itemize}
\end{frame}

\begin{frame}{3.~Dummy-variable regression}
  \begin{defbox}[Seasonal dummies]
    Define $D_{k,t} = 1$ if $t$ is in season $k$, $0$ otherwise.
    \[
      Y_t \;=\; \beta_0 + \sum_{k=2}^{d} \beta_k\, D_{k,t} + \varepsilon_t.
    \]
    Estimated by OLS. Use $d - 1$ dummies (drop one to avoid collinearity).
  \end{defbox}
  \medskip
  \begin{itemize}
    \item Equivalent to seasonal averaging when no other regressor is
          present.
    \item Easy to combine with a polynomial trend in the same regression.
  \end{itemize}
\end{frame}

\begin{frame}{3.~Cosine--sine (Fourier) model}
  \begin{defbox}[Trigonometric seasonality]
    With period $d$ and angular frequency $\omega = 2\pi / d$:
    \[
      S_t \;=\; a\, \cos(\omega t) + b\, \sin(\omega t).
    \]
    Higher harmonics: add $\cos(j \omega t),\, \sin(j \omega t)$ for
    $j = 2, 3, \dots, \lfloor d/2 \rfloor$.
  \end{defbox}
  \medskip
  \begin{itemize}
    \item \textbf{Compact:} 2 parameters per harmonic vs.\ $d - 1$ dummies.
    \item Generalises to non-integer or non-periodic seasonal patterns.
    \item With \emph{all} harmonics, exactly equivalent to dummies
          (linear span argument).
  \end{itemize}
\end{frame}

\begin{frame}[fragile]{3.~Fourier seasonality in Python}
\begin{lstlisting}
import numpy as np
import statsmodels.api as sm

omega = 2 * np.pi / 12
X = np.column_stack([np.cos(omega * t), np.sin(omega * t)])
X = sm.add_constant(X)

model = sm.OLS(y, X).fit()
season = model.predict(X)
\end{lstlisting}
  \medskip
  Add a column \texttt{t} to \texttt{X} for a simultaneous linear trend.
\end{frame}

% =====================================================================
\section{Residuals and diagnostics}

\begin{frame}{4.~Residuals after decomposition}
  \[
    \widehat{\varepsilon}_t \;=\; Y_t - \widehat{T}_t - \widehat{S}_t.
  \]
  \medskip
  \textbf{Quality checks.}
  \begin{itemize}
    \item \textbf{No remaining trend:} plot $\widehat{\varepsilon}_t$
          versus $t$; should fluctuate around $0$.
    \item \textbf{No remaining seasonality:} plot the autocorrelation
          (ACF) — no spike at lag $d$.
    \item \textbf{Constant variance:} no funnel shape; otherwise
          ARCH-type effect (Chapter 6).
    \item \textbf{Approximate normality:} histogram + QQ-plot.
  \end{itemize}
  \medskip
  \begin{infobox}
    The residual $\widehat{\varepsilon}_t$ is the input to ARMA modelling
    in Chapter~5.
  \end{infobox}
\end{frame}

% =====================================================================
\section{Workflow}

\begin{frame}{5.~Full decomposition workflow}
  \begin{enumerate}
    \item Plot $Y_t$. Decide additive vs.\ multiplicative
          (constant or growing seasonal amplitude).
    \item If multiplicative, work on $\log Y_t$.
    \item Estimate $\widehat{T}_t$ (MA, polynomial, or LOESS).
    \item Detrend: $\widetilde{Y}_t = Y_t - \widehat{T}_t$.
    \item Estimate $\widehat{S}_t$ from $\widetilde{Y}_t$.
    \item Residual: $\widehat{\varepsilon}_t = Y_t - \widehat{T}_t - \widehat{S}_t$.
    \item Diagnostics on $\widehat{\varepsilon}_t$.
    \item If residual is stationary: model with ARMA (Chapter 5).
  \end{enumerate}
\end{frame}

\begin{frame}{Up next}
  \textbf{Chapter 3.} Stationary time series.
  \medskip
  \begin{itemize}
    \item Strict and weak stationarity.
    \item Autocovariance function and ACF/PACF.
    \item White noise.
    \item Wold decomposition.
  \end{itemize}
  \medskip
  \begin{center}\Large\bfseries Thank you.\end{center}
  \begin{center}\footnotesize\auteurNom\ -- \auteurEmail\ -- \auteurInstitution\end{center}
\end{frame}

\end{document}
